À $\qty{20}{\degreeCelsius}$ et sous la pression atmosphérique normale, on
dissout $\qty{1.2}{\L}$ de chlorure d'hydrogène dans un volume d'eau tel que
l'on obtienne finalement $\qty{0.50}{\L}$ de solution.
\begin{enumerate}
    \item Calculer la quantité de matière de chlorure d'hydrogène dissous.
    \item Calculer la concentration de la solution obtenue.
    \item Écrire l'équation de dissolution du chlorure d'hydrogène dans l'eau.
          Calculer les concentrations $[\ce{H3O+}]$ et
          $[\ce{Cl-(aq)}]$.
\end{enumerate}
\begin{donnees}
    $V_{m}=\qty{24}{\L\per\mol}$.
\end{donnees}

